\chapter{Geometric Verification}
\label{chap:geometric_verification}
Geometric verification is used to make sure that matched keypoints between images truly correspond to the same physical locations on an animal reducing false positives that arise from visual similarites alone. This step is crucial in refining the initial matches obtained through descriptor-based feature matching by enforcing spatial consistency.

\section{Theoretical Foundation}
Local feature descriptors such as DISK\cite{tyszkiewicz2020disklearninglocalfeatures} or KeyNet combined with AffNet and HardNet \cite{keynet, mishkin2018repeatabilityenoughlearningaffine, Mishchuk2017HardNet}, offer robust initial matches based on visual similarity. However, matches that are generated purely from descriptor similarity may include mistakes, especially in visually cluttered or repetitive environments common in wildlife images. Geometric verification, as suggested in \cite{springerSpeciesAgnosticPatterned}, addresses this by validating matches through spatial coherence checks. 

\section{Establishing Initial Correspondences}
Initial correspondences are established trough descriptor matching using LightGlue\cite{lindenberger2023lightgluelocalfeaturematching}, an attention-based method. Initially methods like Lowe's ratio test were considered but matching with LightGlue proved significantly faster and more accurate, effectively eliminating the need for traditional filtering methods.

\section{Transformation Estimation with RANSAC}
The core of geometric verification involves estimating a homography transformation between matched keypoints using RANSAC (Random Sample Consensus).RANSAC iteratively samples minimal subsets of correspondences, computes candidate transformations and evaluates their agreement with the remaining matches. Matches that fall within a defined spatial threshold are considered inliers. 

\section{Scoring and Verification}
Matches are assessed primarily by their inlier count, the number of matches consistent with the estimated homography, and their spatial distribution:

    \begin{itemize}
        \item \textbf{Inlier count:} Number of correspondences that agree with the estimated transformation; higher counts suggest more reliable matches.
        \item \textbf{Inlier ration: } Normalizes the inlier count by the total number of tentative matches,
  $$
  r = \frac{|I|}{|C|},
  $$
  where $|I|$ denotes the inlier count and $|C|$ the total correspondences.
        \item \textbf{Combined scoring:} Integrates geometric consistency with descriptor similarity (fisher vector distances) to leverage both spatial and appearance information.
    \end{itemize}
Penalties for low inlier counts or poor geometric coherence discourage accepting borderline matches, improving the overall system reliability and accuracy.
